\section{内廷、边地与区域军府的分叉}
\label{sec:late-tang-regional-fracture}

835年甘露之变失败后，宦官杀害朝臣，宫廷中的文官、皇帝与神策军不再拥有可以彼此校正的稳定关系。840年回鹘汗国覆亡，部众向河西、河东迁徙；842年吐蕃赞普遇弑，此后高层冲突扩展；845年武宗废寺并没收部分寺院资产。把这些事合称为“晚唐衰乱”会掩盖它们不同的尺度：宫廷政变关乎禁军与诏令，迁徙关乎边塞供给与安置，吐蕃变局改变走廊政治，废寺则牵涉铜料、户籍、财产和僧俗关系。%
\footnote{甘露之变、回鹘迁徙、吐蕃继承冲突与会昌废佛见新增账本 LT-835-044 至 LT-845-049，依据《旧唐书》文宗、武宗纪、宦官传、回鹘传、吐蕃传及《资治通鉴》。废寺数、还俗数和资产规模来自行政与后世叙事，不能直接当作全国统计。}

张议潮848年驱逐沙州吐蕃守军，851年才获唐廷授归义军节度使；这段间隔说明册命是建立联系的手段，而非朝廷立即恢复河西直辖的证据。安南在863年失守、866年由高骈收复，也同样提示海陆通道、城防和地方军政必须分别维持。敦煌文书能让人看到归义军时期的契约、寺院与军政实践，却不能替开封或长安绘出一张完整的财政地图。%
\footnote{张议潮与归义军、安南失守及收复见 LT-848-051、LT-851-052、LT-863-054、LT-866-055，依据《旧唐书》张议潮、高骈传及《资治通鉴》，并应与敦煌文书和地方考古对读。张议潮所称州数与都护府的行政名号均须同实际控制分开。}

庞勋868年率戍卒起事，王仙芝、黄巢自874年起在关东聚众并逐步分兵，879年黄巢进入广州，880年又攻入长安。广州港市与关中京城的先后易手，不能被写成同一种“流寇”的自然路线：军队要在各地获得粮食、船只、仓储、向导和新成员，地方也会以逃散、交易、抵抗或投附回应。史书尤其难以保存城内普通人的选择，故不能把入城叙事变成全体居民的态度。%
\footnote{庞勋、王仙芝、黄巢起事及广州、长安战事见 LT-868-056、LT-874-058 至 LT-880-062，依据《旧唐书·僖宗本纪》《旧唐书·黄巢传》《新唐书·黄巢传》与《资治通鉴》。兵数、掠获和死亡数多为战时或战后叙事，须避免精确化。}

黄巢政权瓦解后，朱温降唐、沙陀等军收复长安，朝廷反而更依赖能够自行供养和任命将校的军府。885年盐池冲突使僖宗再度出奔；其后王建控制成都、李茂贞巩固凤翔、杨行密获吴王号，显示区域中心并非唐亡才突然出现。900--905年的宫廷废立、诛宦官、迁洛阳与白马驿清洗，则把神策军、朝臣和受禅礼文接入朱温的军政网络。907年唐亡是这条链的结果，却不是区域政治从此失去唐朝名号、仪式和官僚语言的时刻。%
\footnote{朱温降唐至唐亡的关键节点见 LT-882-064 至 LT-907-080，尤其盐池冲突、成都和凤翔军府、吴王册封、宫廷政变与迁洛阳。可互校《旧唐书》昭宗、哀帝纪及朱全忠、王建、李茂贞、杨行密诸传、《新唐书》与《资治通鉴》；后梁和北宋材料都带有事后正统立场。}

\subsection{交叉阅读窗口八：阿房宫的寓言与晚唐的断裂}

甘露之变、会昌废佛与黄巢入长安，先要依《旧唐书》《新唐书》本纪、宦官传和方镇传建立不相混淆的年序，再以《资治通鉴》检验前后连锁。《唐会要》的寺观、军政条目和《通典》的财政分类可见朝廷想如何归类资产、兵员与官署；归义军敦煌文书、张议潮相关碑刻、墓志及吐蕃材料却提示河西的实际重组不等于长安命令的即时延伸。每种材料都带着地点与用途：废寺数、流民数和黄巢军数不应被相加为一套精确的全国衰亡指数。%
\footnote{可参《旧唐书》文宗、武宗、僖宗纪及宦官、张议潮、黄巢诸传，《新唐书》相关纪传、《资治通鉴》唐纪、《唐会要》寺观、军政条目与《通典》食货门。敦煌文书和地方碑刻主要照亮河西个案，不能替代关东和岭南的连续档案。}

杜牧《阿房宫赋》写道：“秦人不暇自哀，而后人哀之；后人哀之而不鉴之，亦使后人而复哀后人也。”这篇大和初年的赋以秦宫的铺张、毁灭和后人回望构成历史寓言；其排比与转折的力量，是把帝国兴亡化为当下读者必须承担的鉴戒。它不是关于835年宫廷政变、845年废寺或875年起事的编年记录，杜牧的士人立场也不能代表宦官、僧侣、军户和商人的判断。正因为如此，文学文本适合打断“唐亡必然”的舒适叙事，而不适合填充事件数字。它要求我们回到多种材料，分别追问宫廷军权、区域财政和跨境迁徙怎样在不同地点断裂。
